% !TEX root = ../main.tex
\cleardoublepage
\chapter*{Abstract - Liam}
\markboth{Abstract}{Abstract}
\addcontentsline{toc}{chapter}{Abstract}

This diploma thesis details the development of ``Ascenta'', an intelligent assistive wearable designed to support visually impaired individuals. Given the challenges of navigating complex urban environments, the project aims to create an affordable solution using artificial intelligence to provide information about the user's social surroundings.

The technical framework utilizes the Arduino Nicla Vision microcontroller. The firmware implements the \textbf{F}aster \textbf{O}bjects, \textbf{M}ore \textbf{O}bjects (FOMO) architecture for efficient on-device face detection. For user interaction, a local and privacy-oriented voice recognition system was developed using the \textbf{Vosk engine} within a custom Android application. Communication is established via an optimized TCP protocol to ensure low-latency data streaming between the wearable and the mobile device.

The current implementation successfully recognizes faces within the field of view and determines proximity using Time-of-Flight sensors. The results demonstrate the efficiency of the decentralized architecture and provide a solid foundation for future development of general obstacle detection models.

\cleardoublepage
\chapter*{Kurzfassung - Liam}
\markboth{Kurzfassung}{Kurzfassung}
\addcontentsline{toc}{chapter}{Kurzfassung}

Die vorliegende Diplomarbeit dokumentiert die Entwicklung von \glqq Ascenta\grqq, einer intelligenten Assistenzbrille zur Unterstützung sehbeeinträchtigter Personen im Alltag. Angesichts der Komplexität moderner urbaner Umgebungen besteht das Ziel des Projekts in der Realisierung eines kostengünstigen Wearables, das mittels künstlicher Intelligenz Informationen über das soziale Umfeld liefert.

Die technische Umsetzung basiert auf dem Arduino Nicla Vision Mikrocontroller. Die Firmware nutzt die FOMO-Architektur (Faster Objects, More Objects) zur effizienten On-Device-Gesichtserkennung. Zur Interaktion wurde eine lokale, datenschutzkonforme Sprachverarbeitung auf Basis der \textbf{Vosk-Engine} in einer Android-Applikation implementiert. Die Kommunikation erfolgt über ein optimiertes TCP-Protokoll, um minimale Latenzzeiten beim Datenstreaming zwischen der Brille und dem Smartphone zu gewährleisten.

Das realisierte System ist in der Lage, Personen im Sichtfeld zu identifizieren und die Distanz mittels Time-of-Flight-Sensorik zu bestimmen. Die Ergebnisse belegen die Funktionalität der dezentralen Architektur und bilden die Grundlage für zukünftige Erweiterungen zur allgemeinen Hinderniserkennung.